%%%%%%%%%%%%%%%%%%%%%%%%%%%%%%%%%%%%%%%
% Cover Letter - Rendezvous Robotics
% Summer 2027 Software Engineering Intern, Golden CO
% Plain XeLaTeX — matches simple resume style (Spectral font)
%%%%%%%%%%%%%%%%%%%%%%%%%%%%%%%%%%%%%%%

\documentclass[11pt]{article}

\usepackage[margin=0.75in, top=0.7in, bottom=0.7in]{geometry}
\usepackage[hidelinks]{hyperref}
\usepackage{parskip}
\usepackage{fontspec}

\setmainfont{Spectral}[
  Path = ../../cv/fonts/,
  UprightFont    = Spectral-Regular.ttf,
  BoldFont       = Spectral-Bold.ttf,
  ItalicFont     = Spectral-Italic.ttf,
  BoldItalicFont = Spectral-BoldItalic.ttf
]

\setlength{\parindent}{0pt}
\setlength{\parskip}{14pt}
\linespread{1.15}
\raggedright

% No hyphenation
\hyphenpenalty=10000
\exhyphenpenalty=10000

\pagestyle{empty}
\urlstyle{same}

\begin{document}

% ── DATE ─────────────────────────────────────────────────────────────────────
\begin{flushright}
  \today
\end{flushright}

% ── BODY ─────────────────────────────────────────────────────────────────────
Dear Rendezvous Robotics Hiring Team,

I'm applying for the Summer 2027 Software Engineering Internship. I'm a
Computer Science student at RIT in the combined bachelor's and master's
program, with an expected graduation date of December 2027. At GE Vernova,
I built a full I/O subsystem from scratch in C on an STM32F7 for
industrial radios. On the RIT Electric Vehicle Team, I'm the vehicle
control unit (VCU) Software Lead, building real-time ThreadX firmware for
our track motorcycle. In my free time I'm building an autonomous
fixed-wing flight controller on an ESP32S3, with PID stabilization, sensor
fusion, and GPS waypoint following validated across an 80-test SITL
regression suite.

What got me interested in Rendezvous Robotics is that the tiles physically
dock and reconfigure into something bigger than what fits inside a
rocket's fairing. The part of the posting that stood out most is keeping
communication between units working as the assembly itself changes shape
or membership. On the VCU, I designed the CAN and CANopen message formats
connecting it to the motor controller, battery management system,
handlebar interface board, and low-voltage subsystem, and had to handle
cases where a node drops off the bus or comes back online without taking
the rest of the vehicle down with it. That's a smaller version of the same
problem, a fixed set of devices with known roles. A flight assembly where
units can join or leave entirely changes what the protocol has to
guarantee, and figuring out how to handle that is genuinely interesting to
me.

I've also built the kind of tooling this role describes for ground
station operations. At GE Vernova, I built a Python/Flask dashboard that
polls and visualizes Modbus diagnostic data in real time, with historical
trend data, as an internal engineering tool. I
build things in simulation before I trust them on hardware, too. My
flight controller project's SITL simulators for attitude and navigation
logic are the clearest example.

Real-time embedded C is where I spend most of my time. I've structured the
VCU's ThreadX design with a timer-triggered control thread, a health
monitoring thread, and interrupt-fed CAN queues protected by mutexes and
event flags, and I debug at the hardware level regularly with a Saleae
logic analyzer, an oscilloscope, and a J-Link Segger, including tracing a
servo fault at GE Vernova back to a bad servo rather than the signal
driving it.

I would welcome the opportunity to talk more about where I could
contribute. Thank you for your consideration.

Thanks,

Travis Brown

\end{document}
